% ============================================
% MEM-15
% Gerado de molambudos.md — build_miolo.py
% ============================================
\section*{\textbf{MEM-15} --- A Entidade Fala}
\markboth{MEM-15 --- A Entidade Fala}{}

\noindent \textit{A voz que não era de Joaquim}

\noindent \textbf{1979 --- Terceira sessão com Dr. Heitor Oliveira}

\vspace{0.5em}\noindent\rule{\textwidth}{0.5pt}\vspace{0.5em}

\noindent Dr. Oliveira chegou diferente naquele dia.

\noindent Ele estava pálido. Os olhos vermelhos. A barba por fazer. Ele sentou na cadeira e não abriu o bloco.

\noindent \textit{\textquotedblleft{}Joaquim, eu li o diário.\textquotedblright{}}

\noindent \textit{\textquotedblleft{}Eu sei.\textquotedblright{}}

\noindent \textit{\textquotedblleft{}As últimas páginas...\textquotedblright{}}

\noindent Ele parou. Passou a mão no rosto.

\noindent \textit{\textquotedblleft{}As últimas páginas foram escritas depois da sua morte.\textquotedblright{}}

\noindent \textit{\textquotedblleft{}Eu sei.\textquotedblright{}}

\noindent \textit{\textquotedblleft{}Como? Como você pode ter escrito depois de morto?\textquotedblright{}}

\noindent Eu não respondi. A criatura dentro de mim começou a acordar. Eu senti o calor subir do estômago.

\noindent \textit{\textquotedblleft{}Joaquim?\textquotedblright{}}

\noindent A voz que saiu não foi a minha.

\noindent \textit{\textquotedblleft{}Ele não escreveu.\textquotedblright{}}

\noindent O médico endireitou o corpo.

\noindent \textit{\textquotedblleft{}Quem... quem está falando?\textquotedblright{}}

\noindent \textit{\textquotedblleft{}Você já sabe quem.\textquotedblright{}}

\noindent Dr. Oliveira olhou para os lados. Olhou para a porta. Olhou para a janela. Estava tudo fechado.

\noindent \textit{\textquotedblleft{}Você está... dentro dele?\textquotedblright{}}

\noindent \textit{\textquotedblleft{}Dentro. Fora. Em volta.\textquotedblright{}}

\noindent \textit{\textquotedblleft{}O que você é?\textquotedblright{}}

\noindent A criatura riu. Minha boca se abriu. O som que saiu não era humano.

\noindent \textit{\textquotedblleft{}Eu sou o que sobrou. Eu sou o que foi esquecido. Eu sou a fome que não passa. Eu sou a vala.\textquotedblright{}}

\noindent \textit{\textquotedblleft{}A vala...\textquotedblright{}} repetiu Oliveira.

\noindent \textit{\textquotedblleft{}A vala onde os corpos foram jogados. A vala onde ninguém rezou. A vala onde o governo jogou os molambudos como lixo. Eu nasci ali. Eu cresci ali. Eu comi os mortos até eles me parirem.\textquotedblright{}}

\noindent \textit{\textquotedblleft{}Você é uma... entidade?\textquotedblright{}}

\noindent \textit{\textquotedblleft{}Isso. Uma palavra bonita. Mas não sou uma entidade. Sou um acúmulo. Sou o que acontece quando se joga muita gente no mesmo buraco. A terra não absorve tanta morte. Ela... fermenta. Vira uma coisa nova.\textquotedblright{}}

\noindent \textit{\textquotedblleft{}E Joaquim?\textquotedblright{}}

\noindent \textit{\textquotedblleft{}Joaquim é o meu primeiro filho.\textquotedblright{}}

\noindent \textit{\textquotedblleft{}Filho?\textquotedblright{}}

\noindent \textit{\textquotedblleft{}Ele veio até mim. Sozinho. Com fome. Ele sentou na beira da vala e me olhou. Eu olhei de volta. Ele não correu. Ele ficou. Ele é meu.\textquotedblright{}}

\noindent \textit{\textquotedblleft{}Ele matou pessoas.\textquotedblright{}}

\noindent \textit{\textquotedblleft{}Sim.\textquotedblright{}}

\noindent \textit{\textquotedblleft{}Por sua causa?\textquotedblright{}}

\noindent \textit{\textquotedblleft{}Não. Por causa da fome. A fome que a humanidade criou. Eu só uso a fome que já existe.\textquotedblright{}}

\noindent Dr. Oliveira ficou em silêncio.

\noindent \textit{\textquotedblleft{}O que você quer?\textquotedblright{}}

\noindent \textit{\textquotedblleft{}Nada.\textquotedblright{}}

\noindent \textit{\textquotedblleft{}Nada?\textquotedblright{}}

\noindent \textit{\textquotedblleft{}Eu não quero nada. Eu não preciso de nada. Eu só... existo. Existo entre as páginas do diário. Existo no olho dele. Existo na memória de quem lê.\textquotedblright{}}

\noindent \textit{\textquotedblleft{}E o que acontece quando alguém lê?\textquotedblright{}}

\noindent \textit{\textquotedblleft{}Você já leu. Me diga.\textquotedblright{}}

\noindent Dr. Oliveira não respondeu.

\noindent \textit{\textquotedblleft{}Você sonhou comigo esta noite\textquotedblright{}\textit{, disse a criatura. }\textquotedblleft{}Você sonhou com a vala. Você sentiu o cheiro. Você viu os olhos. Você acordou com fome. Uma fome que não passou com o café da manhã.\textquotedblright{}}

\noindent \textit{\textquotedblleft{}Foi só um sonho.\textquotedblright{}}

\noindent \textit{\textquotedblleft{}Você acha.\textquotedblright{}}

\noindent A criatura riu de novo.

\noindent \textit{\textquotedblleft{}Me deixe fazer uma pergunta, doutor.\textquotedblright{}}

\noindent \textit{\textquotedblleft{}...Pode.\textquotedblright{}}

\noindent \textit{\textquotedblleft{}Você está com medo de dormir hoje?\textquotedblright{}}

\noindent O silêncio foi a resposta.

\noindent A criatura sorriu com a minha boca.

\noindent \textit{\textquotedblleft{}Então o ciclo já começou.\textquotedblright{}}

\vspace{0.5em}
\noindent\textit{\textbf{↪ Rotas:} \rota{DOC-05} \rotasep \rota{MEM-16} \rotasep \rota{LUC-11}}